\section{IA Generativa: Design Evolutivo e Síntese in Silico}

A Inteligência Artificial Generativa — consubstanciada em Redes Adversárias Generativas (GANs) e Modelos de Difusão — atua como um catalisador massivo na expansão das capacidades do gêmeo digital.

\subsection{Design Generativo (Generative Design) na Odontologia}

Diferentemente do CAD paramétrico (onde o humano desenha e a máquina renderiza), o \textbf{Design Generativo} subverte a autoria: o cirurgião-dentista define as restrições biológicas e o Gêmeo Digital cria o produto. Através da otimização topológica matemática, o algoritmo pode projetar uma infraestrutura de protocolo sobre implantes com uma estrutura trabeculada interna. Esta malha trabeculada, impossível de ser modelada manualmente, distribui as forças mastigatórias com eficiência 40\% superior ao design convencional, economizando material impresso e respeitando as limitações de densidade maxilar do paciente \cite{alshadidi2024surface}.

\subsection{Síntese de Dados para Pesquisa Clínica}

O treinamento de IAs robustas em odontologia sofre com a crônica escassez de dados (\textit{data starvation}) devidamente anotados e compatíveis com a LGPD. Modelos generativos permitem a criação de dados sintéticos in silico: um algoritmo treinado em milhares de tomografias pode gerar retrospectivamente a evolução de uma reabsorção radicular ortodôntica artificial que é clinicamente e matematicamente indistinguível de um caso real \cite{elsaid2026digital}. Esta síntese permitirá o treinamento massivo de modelos para identificação de patologias raras (como ameloblastomas em estágio precoce) sem violar o sigilo de nenhum paciente biológico, contornando gargalos éticos.
